% titel: Blatt 9, A1: Konzentrationsmasse einer Mischung
% skript: Mortimer Kap. 13 (Lösungen, Konzentrationsmasse)

% Teil I: Werkzeuge, kein Lösungsweg.

Zum Beschreiben einer Lösung gibt es mehrere Grössen, die sich nur im
Nenner unterscheiden.

\begin{werkzeugkasten}[Konzentrationsmasse]
Allgemein braucht man die Stoffmengen $n_i$ der Komponenten:
\[
  n = \frac{m}{M}\qquad\text{(bzw. } n=\frac{\rho V}{M}\text{)}.
\]
\begin{itemize}
  \item \textbf{Stoffmengenkonzentration c} = Anzahl der Mol gelöster
        Substanz pro Liter \emph{Lösung}:
        \[
          c = \frac{n_{\text{gelöst}}}{V_{\text{Lösung}}}.
        \]
  \item \textbf{Molalität b} = Mol gelöster Substanz pro Kilogramm
        \emph{Lösungsmittel}:
        \[
          b = \frac{n_{\text{gelöst}}}{m_{\text{Lösungsmittel}[kg]}}.
        \]
  \item \textbf{Molenbruch x} = Molanteil an der Gesamtzahl aller Mol:
        \[
          x_i = \frac{n_i}{\sum_j n_j}.
        \]
\end{itemize}
Volumen der Lösung aus Masse und Dichte: $V = m_{\text{ges}}/\rho$.
\textit{Skript:} Mortimer Kap. 13 (Konzentrationen).
\end{werkzeugkasten}

\begin{achtung}[Darauf achten]
\begin{itemize}
  \item Bei c steht die \emph{Lösung} im Nenner (aus $\rho$ und Gesamtmasse
        berechnen), bei b nur das \emph{Lösungsmittel} und in kg.
  \item Molenbruch aller Komponenten summiert sich zu 1.
  \item Wasser wird hier als Lösungsmittel betrachtet; die anderen beiden
        Stoffe sind gelöst.
  \item Erst alle $n_i$ ausrechnen, dann einsetzen.
\end{itemize}
\end{achtung}
